\documentclass[11pt]{article}

% Change "review" to "final" to generate the final (sometimes called camera-ready) version.
% Change to "preprint" to generate a non-anonymous version with page numbers.
\usepackage[preprint]{acl}

% Standard package includes
\usepackage{times}
\usepackage{latexsym}

% For proper rendering and hyphenation of words containing Latin characters (including in bib files)
\usepackage[T1]{fontenc}
% For Vietnamese characters
% \usepackage[T5]{fontenc}
% See https://www.latex-project.org/help/documentation/encguide.pdf for other character sets

% This assumes your files are encoded as UTF8
\usepackage[utf8]{inputenc}
\usepackage{parskip}

% This is not strictly necessary, and may be commented out,
% but it will improve the layout of the manuscript,
% and will typically save some space.
\usepackage{microtype}

% This is also not strictly necessary, and may be commented out.
% However, it will improve the aesthetics of text in
% the typewriter font.
\usepackage{inconsolata}

%Including images in your LaTeX document requires adding
%additional package(s)
\usepackage{graphicx}

% If the title and author information does not fit in the area allocated, uncomment the following
%
%\setlength\titlebox{<dim>}
%
% and set <dim> to something 5cm or larger.

\title{From MIDI to track generation using NLP Techniques}

% Author information can be set in various styles:
% For several authors from the same institution:
% \author{Author 1 \and ... \and Author n \\
%         Address line \\ ... \\ Address line}
% if the names do not fit well on one line use
%         Author 1 \\ {\bf Author 2} \\ ... \\ {\bf Author n} \\
% For authors from different institutions:
% \author{Author 1 \\ Address line \\  ... \\ Address line
%         \And  ... \And
%         Author n \\ Address line \\ ... \\ Address line}
% To start a separate ``row'' of authors use \AND, as in
% \author{Author 1 \\ Address line \\  ... \\ Address line
%         \AND
%         Author 2 \\ Address line \\ ... \\ Address line \And
%         Author 3 \\ Address line \\ ... \\ Address line}

\author{Thierry Poey \\ 
\texttt{thierry.poey@umontreal.ca} \\\And
  Michael Xi \\ 
  \texttt{michael.xi@umontreal.ca} \\\And
  Rabbi Binda Mvitu \\
  \texttt{rabbi.binda.mvitu@umontreal.ca} 
  }

%\author{
%  \textbf{First Author\textsuperscript{1}},
%  \textbf{Second Author\textsuperscript{1,2}},
%  \textbf{Third T. Author\textsuperscript{1}},
%  \textbf{Fourth Author\textsuperscript{1}},
%\\
%  \textbf{Fifth Author\textsuperscript{1,2}},
%  \textbf{Sixth Author\textsuperscript{1}},
%  \textbf{Seventh Author\textsuperscript{1}},
%  \textbf{Eighth Author \textsuperscript{1,2,3,4}},
%\\
%  \textbf{Ninth Author\textsuperscript{1}},
%  \textbf{Tenth Author\textsuperscript{1}},
%  \textbf{Eleventh E. Author\textsuperscript{1,2,3,4,5}},
%  \textbf{Twelfth Author\textsuperscript{1}},
%\\
%  \textbf{Thirteenth Author\textsuperscript{3}},
%  \textbf{Fourteenth F. Author\textsuperscript{2,4}},
%  \textbf{Fifteenth Author\textsuperscript{1}},
%  \textbf{Sixteenth Author\textsuperscript{1}},
%\\
%  \textbf{Seventeenth S. Author\textsuperscript{4,5}},
%  \textbf{Eighteenth Author\textsuperscript{3,4}},
%  \textbf{Nineteenth N. Author\textsuperscript{2,5}},
%  \textbf{Twentieth Author\textsuperscript{1}}
%\\
%\\
%  \textsuperscript{1}Affiliation 1,
%  \textsuperscript{2}Affiliation 2,
%  \textsuperscript{3}Affiliation 3,
%  \textsuperscript{4}Affiliation 4,
%  \textsuperscript{5}Affiliation 5
%\\
%  \small{
%    \textbf{Correspondence:} \href{mailto:email@domain}{email@domain}
%  }
%}

\begin{document}
\maketitle

\begin{abstract}
This project studies music generation from MIDI files by adding simple NLP module to existing open-source pipelines. The main goal is not to build a new SOTA model but to evaluate whether existing models can be adapted to an end-to-end format from a user prompt to a usable musical outputs. More specifically, we study whether a language model can influence MIDI generation from a simple text prompt.

In this first phase, we reviewed recent symbolic music generation system such as MIDI-GPT and other text to music pipelines. We then built a simple NLP module where an LLM converts a text prompt into a small DSL representation, which is parsed and cleaned to match the input parameters expected by MIDI-GPT. This setup is simple and is mainly used to observe the effect of the prompt and the limits of this integration.

Our preliminary results are mostly based on qualitative listening and comparisons between original MIDI files and generated outputs. The first observations show that the prompt has influence but this influence is often ambiguous, and hard to keep over long sequences. We also found important issues with reproducibility, documentation, and experimental cost in the open source models we study. A good conclusion is that these results suggest that post-processing may be a more realistic direction than building a fully new MIDI generation pipeline.
\end{abstract}

\section{Introduction}

With the rise of language model, music generation from MIDI data is an active research topic. These methods make it possible to interact with music generation systems through high-level prompts and to design end-to-end pipelines that go from user input to usable music output. However, these systems remain difficult to reuse locally and reproducible open source setting. Another major issue is the alignment between the user request and the generated music.

In this first phase, we explore a simple NLP module to guide the generation of an existing MIDI model. Indeed, we use a prompt sent to an LLM, convert it into an intermediate representation and transform it into parameters that can be used by MIDI-GPT. The goal of this step is not yet to build a complete musical system but to understand how far simple text conditioning can influence generation and what the main limits of this approach are.

Our exploration discover several difficulties: missing transparency about the data and training pipelines, very weak documentation and code not commented , and high computational cost even for a small models. In addition, our first qualitative evaluations show that differences between generations with and without prompts can sometimes be heard. For example, simple oppositions such as happy versus sad remain fragile and often unclear particularly for long sequences.

\section{Related work}

Symbolic music generation from MIDI files has been studied in recent years. Some approaches model note sequences or MIDI events in an autoregressive way while others use text conditioning to guide the generation process.

Among the models closest to our work, MIDI-GPT \cite{Pasquier2025MIDIGPT} is an important reference because it generates MIDI sequences with a GPT architecture. Other models such as MIDI-LLM \cite{Wu2025MIDILLM} and Text2Midi \cite{Bhandari2025Text2midi} explore the link between text descriptions and MIDI generation. Some models don't focus on generation but on task such as inpainting, continuation or partial rewriting of musical sequences.

\section{Methods}

Our approach adds an LLM prompt step before an existing MIDI generation model. For now, we use a ollama model to transform a text prompt into a generic JSON DSL. This DSL contains different control fields that can later be refined. The goal is to build a first generic version that the model can use easily create before design the decision more precisely.

After this step, the output must be cleaned before it can be used. First, we normalize the values to remove issues such as white spaces, uppercase letters or bad characters. We also rewrite synonyms. This step is a form of canonicalization. Then, we validate the output so that it matches an allowed schema.

We also need to interpret some values because the LLM does not directly produce an exact representation. For example, a mood such as calm or happy must be translated into more concrete controls. A calm mood may correspond to lower density and less variation while an energetic mood may correspond to higher density. These choices does not constitute the truth because it's a creative domain but give a basic direction, in particular because we are limited by the parameters in MIDI-GPT.

After cleaning and interpretation, the DSL is converted into the parameter names expected by MIDI-GPT. For example, a low polyphony level is mapped to the correct internal polyphony settings (based on the model norms). In addition, a MIDI input file is used to indicate which bars should be inpainted during generation.

Instead of retraining a large model or building a new one, we focus on the effect of the prompt.

\section{Dataset}

For the moment, we didn't build or train on a new dataset because open source models such as MIDI-GPT already provide pretrained checkpoints. Our work focuses on reusing these existing resources and evaluating them locally. MIDI-GPT is trained on the GigaMIDI dataset which contains about 2.1 million MIDI files divided into three subsets: all-instruments-with-drums, drums-only, and no-drums.

We also considered using the MIDICaps dataset which provides pairs of MIDI files + text captions. This dataset could be useful for comparing the user input with the generated MIDI output. It is also the dataset used to train Text2midi.

The data used in this project are MIDI files. This format represents musical events such as notes, duration, pitch, velocity and timing. In our case, MIDI files are both the source used for generation and a reference for comparing outputs. However, the documentation is often too limited. In several cases, the data description and preprocessing steps are not clearly explained. This makes reproduction difficult and also raises questions about \textit{building a pipeline with our own generation model}.

\section{Baseline and evaluation}

The main baseline in our project is the original MIDI-GPT. This baseline was obtained from the existing MIDI-GPT codebase and adapted to our local setup. We use this baseline because it allows a direct comparison between the original model behavior and our prompt version. 

Our evaluation is mainly listening. The qualitative evaluation is based on comparisons between generations with and without prompt conditioning, between short and long generations, and between simple contrasting prompts such as happy and sad. In the next stage, we plan to make this protocol more precise by using clear comparison criteria.

\section{Experiment details}

The experiments were run locally on an RTX 5070 Ti. This directly affect the choice of the model, the size of the data subsets that could be used and the number of experiments we could run. In particular, training on a small model such as GPT2 small can require 3 days of computation on only a subset of data. Because of this, full retraining is not practical in our setting.

The prompt is first converted into a small JSON DSL with eight variables: instrument, mood, complexity, polyphony\_level, density\_level, register, note\_duration\_level and genre. A post-processing step then normalizes these values, applies default values when needed, provide synonyms and resolves simple dependencies between variables such as links between mood, complexity, density, and register.

The input MIDI file is then converted into JSON with the MIDI-GPT encoder, and only the first track is currently used as the target for generation. Three generation modes are tested: tail\_infill which keeps the first bar as context and generates the rest, full\_conditional which generates all bars and autoregressive which generates the sequence step by step.

The parameters controls are then converted into parameters compatible with MIDI-GPT. The generation temperature is also adjusted based on mood and complexity. The main hyperparameters currently used are batch\_size = 1, max\_steps = 200, polyphony\_hard\_limit = 6, and max\_attempts = 3. Finally, the generated output is converted back into a MIDI file. 

\section{Results}



\section{Future work}

The main problem is that we cannot build a new generative model, so we can only modify either the input or the output. The input seems to be very sensitive and hard to control. Because of this, we prefer to keep the input more general and focus on post-processing techniques. For the final stage, we want to focus first on improving the musical quality of the generated outputs to produce results that sound more musical and less robotic. This includes improving transitions between generated parts, improving the global coherence of the piece and reducing artificial effects that can appear in the outputs. We also want to move toward a more modular system. One limitation of the current approach is that a small change in the prompt can affect the whole generated result. In the next phase, we would like to study whether some properties of the output can be modified more directly, so that the control becomes more local and more stable.

\section{Task Division}

\bibliography{custom}

\end{document}
